% ==================================================================
% id:       comb-018
% topic:    קומבינטוריקה
% subtopic: המקדם הבינומי
% solution: full
% ==================================================================
בבניין 14 קומות. בכל קומה גרות 4 משפחות, חוץ מקומה 4, שבה גרות רק 2 משפחות. בוחרים 4 משפחות מהבניין לוועד הבית.
\begin{enumerate}
\item בכמה דרכים אפשר לבחור את הוועד?
\begin{solbox}
בבניין $13\cdot4+2=54$ משפחות. בוחרים 4 מתוכן, בלי סדר ובלי חזרות: $\binom{54}4=\ans{316{,}251}$.
\end{solbox}
\item בכמה דרכים אפשר לבחור את הוועד כך שכל אחת מהקומות 1, 2, 3 ו-4 תהיה מיוצגת?
\begin{solbox}
בוועד 4 משפחות, ולכן בוחרים בדיוק משפחה אחת מכל אחת מהקומות 1--4. בקומות 1--3 יש 4 אפשרויות בכל קומה, ובקומה 4 יש 2: $4\cdot4\cdot4\cdot2=\ans{128}$.
\end{solbox}
\item בכמה דרכים אפשר לבחור את הוועד כך ששתי המשפחות מקומה 4 יהיו בו?
\begin{solbox}
שתי המשפחות מקומה 4 בוועד, ונשאר לבחור 2 משפחות מתוך 52 המשפחות האחרות: $\binom{52}2=\ans{1326}$.
\end{solbox}
\item בכמה דרכים אפשר לבחור את הוועד כך שקומה 4 תהיה מיוצגת?
\hint{קל יותר לספור את המשלים: בכמה דרכים קומה 4 \emph{לא} מיוצגת?}
\begin{solbox}
נספור דרך המשלים. מספר הוועדים שאין בהם משפחה מקומה 4 הוא $\binom{52}4=270{,}725$. לפי עקרון החיסור, מספר הוועדים שקומה 4 מיוצגת בהם הוא $316{,}251-270{,}725=\ans{45{,}526}$.

\textbf{דרך נוספת.} נפריד לשני מקרים זרים. משפחה אחת מקומה 4: $\binom21\binom{52}3=2\cdot22{,}100=44{,}200$. שתי המשפחות מקומה 4: $1326$ (סעיף ג). לפי עקרון החיבור: $44{,}200+1326=45{,}526$.
\end{solbox}
\end{enumerate}
